% !PRINTED 579
\weberpage{579}
%<<CONT
die Koeffizienten von $A(x)$ ausdrückbar, ist also eine ganze rationale
Funktion von $\lambda$ mit rationalen Zahlenkoeffizienten von der Form:
\weq{12}{S = S_0\lambda^s + S_1\lambda^{s-1} + S_2\lambda^{s-2} + \cdots + S_s,}
worin die $S_0, S_1, S_2, \ldots$, rationale Zahlen sind. Daß es \gesperrt{ganze
Zahlen} sind, soll eben bewiesen werden. Das ergibt sich aber
sehr einfach, wenn man in (12) für $\lambda$ die Entwickelung (3) einsetzt
 und dann die Koeffizienten gleich hoher Potenzen von $q$ auf
beiden Seiten miteinander vergleicht. Man bekommt nämlich zur
Bestimmung der $S_0, S_1, S_2, \ldots$, eine Reihe linearer Gleichungen,
deren jede folgende nur ein neues $S_i$ und zwar mit dem Koeffizienten
 1 enthält, während auf der anderen Seite dieser Gleichungen
 ganze algebraische Zahlen stehen. Hiernach sind also
die $S_0, S_1, S_2, \ldots$, ganze algebraische Zahlen, und da sie rational
sind, sind es auch \gesperrt{ganze rationale Zahlen}.

Da wir überdies nach § 57, II und (7) $A(0)$ und $A(\infty)$
kennen, so ergibt sich für $A(x)$ der Ausdruck:
\weq{13}{A(x) = \pm x^{m^2-1} + a_1 x^{m^2-3} + a_2 x^{m^2-5} + \cdots + m,}
\gesperrt{worin die $a_1, a_2, \ldots$, ganze ganzzahlige Funktionen von
$\lambda$ sind.}

Ferner ist nach § 57, (7):
\[
D(x) = \pm x^{m^2-1} A\left(\frac{1}{x}\right),
\]
und folglich:
\weq{14}{\pm D(x) = \pm 1 + a_1 x^2 + a_2 x^4 + \cdots + m x^{m^2-2}.}
\[
\text{[Das Zeichen } \pm \text{ in (13) und (14) ist nach § 57 bestimmt durch}
(-1)^{\frac{m-1}{2}}\text{].}
\]
Für die Funktionen $B$ und $C$ gilt die Relation:
\weq{15}{B(x) = x^{m^2-1} C\left(\frac{1}{x}\right),}
und ihre Wurzeln sind daher zueinander reziprok. Der erste und
der letzte Koeffizient sind in $B$ und in $C$ gleich der Einheit.

Die Wurzeln von $B$ sind:
\[
y_{h,h'} = \sqrt{\varkappa}\,\mathrm{sn}(\Omega_{h,h'}+K) = \sqrt{\varkappa}\,\frac{\mathrm{cn}\,\Omega_{h,h'}}{\mathrm{dn}\,\Omega_{h,h'}}
\]
\weq{16}{= \frac{\vartheta_{10}\left(\frac{h+h'\omega}{m},\omega\right)}{\vartheta_{00}\left(\frac{h+h'\omega}{m},\omega\right)}}
%CONT>>
